\documentclass[11pt,a4paper]{article}
\usepackage[utf8]{inputenc}
\usepackage[T1]{fontenc}
\usepackage{lmodern}
\usepackage{geometry}
\geometry{margin=1in}
\usepackage{hyperref}
\usepackage{enumitem}
\usepackage{xcolor}
\usepackage{titlesec}
\usepackage{tcolorbox}

\definecolor{primary}{RGB}{33, 64, 154}
\titleformat{\section}{\Large\bfseries\color{primary}}{}{0em}{}[\titlerule]
\titleformat{\subsection}{\large\bfseries\color{darkgray}}{}{0em}{}

\title{\textbf{\Huge LedgerPilotAI} \\ \vspace{0.5em} \Large Step-by-Step Implementation Plan}
\author{Based on the AI Finance Controller Technical Design}
\date{}

\begin{document}

\maketitle
\thispagestyle{empty}

\begin{abstract}
\noindent This document provides a structured, production-oriented implementation roadmap for \textbf{LedgerPilotAI}. The system is designed to automate financial reconciliation, manage exceptions, provide AI-assisted matching, and analyze cash positions across multiple financial data sources.
\end{abstract}

\vspace{1em}

\section{Phase 1: Foundation}
\textit{Objective: Establish the core application infrastructure and database schema.}
\begin{itemize}[label=$\bullet$, leftmargin=*]
    \item \textbf{Frontend:} Initialize the React application with TypeScript, Vite, React Router, and Tailwind CSS.
    \item \textbf{Backend API:} Initialize the Node.js and Express API using TypeScript. Set up the foundational folder structure (routes, controllers, services).
    \item \textbf{Database:} Provision a Neon PostgreSQL database.
    \item \textbf{ORM Setup:} Configure Drizzle ORM, define the core schema (Tenants, Users, Transactions, Exceptions), and implement initial database migrations.
\end{itemize}

\section{Phase 2: Data Pipeline}
\textit{Objective: Enable secure and idempotent ingestion of financial records.}
\begin{itemize}[label=$\bullet$, leftmargin=*]
    \item \textbf{Synthetic Data:} Build a script to generate a synthetic dataset of 100--120 records containing realistic imperfections (e.g., date offsets, amount mismatches, missing references).
    \item \textbf{Ingestion API:} Implement CSV and JSON file upload and ingestion pipelines.
    \item \textbf{Validation:} Add application-level schema validation using Zod (e.g., required external IDs, numeric amounts).
    \item \textbf{Normalization:} Convert incoming records into a canonical transaction format (normalizing names, dates, and minor currency units).
    \item \textbf{Idempotency:} Implement duplicate fingerprinting (SHA256 hash of source, ID, amount, date) to prevent double processing.
\end{itemize}

\section{Phase 3: Reconciliation Engine}
\textit{Objective: Implement the deterministic matching logic and scoring algorithms.}
\begin{itemize}[label=$\bullet$, leftmargin=*]
    \item \textbf{Layer 1 (Exact):} Implement exact matching for references, identifiers, and exact amounts.
    \item \textbf{Candidate Generation:} Build broad database filters to group potential transaction candidates efficiently (e.g., matching currency, dates within a tight tolerance).
    \item \textbf{Deterministic Scoring:} Implement component scores (Amount Score, Date Score, Reference Score, Description Score) and a weighted composite score.
    \item \textbf{Run Model:} Create the \texttt{reconciliation\_runs} logic to track batch progress, match rates, and unresolved amounts.
    \item \textbf{Results Persistence:} Save high-confidence matched pairs to the database and queue low-confidence items as exceptions.
\end{itemize}

\section{Phase 4: AI Layer}
\textit{Objective: Introduce AI safely to handle ambiguous transactions and natural language queries.}
\begin{itemize}[label=$\bullet$, leftmargin=*]
    \item \textbf{Structured Output:} Define strict TypeScript interfaces and JSON schemas for the AI agent's decision-making output (MATCH, REVIEW, UNMATCHED).
    \item \textbf{Tool Calling:} Build validated backend tools that the LLM can invoke (e.g., \texttt{getTransactions()}, \texttt{findCandidates()}).
    \item \textbf{Ambiguous Workflow:} Route medium-confidence fuzzy matches to the AI for semantic evaluation, extracting evidence and reason codes.
    \item \textbf{Explanation Engine:} Use the AI to generate readable explanations for unresolved exceptions.
    \item \textbf{Natural Language Interface:} Implement the AI Finance Assistant endpoint to answer queries (e.g., "Why are there 12 exceptions?").
\end{itemize}

\section{Phase 5: Dashboard \& UI}
\textit{Objective: Build the human-in-the-loop interfaces for finance operators.}
\begin{itemize}[label=$\bullet$, leftmargin=*]
    \item \textbf{Main Dashboard:} Display aggregate metrics, total records processed, match rates, and active exceptions.
    \item \textbf{Run Details Screen:} Build a view to track specific reconciliation runs and their resulting transaction states.
    \item \textbf{Exception Review UI:} Create a queue where human operators can inspect candidate evidence, approve/reject AI suggestions, and add resolution notes.
    \item \textbf{Cash Position View:} Implement the "Why Is My Cash Off?" screen, calculating expected vs. actual balance and tracing the variance to specific unreconciled transactions.
\end{itemize}

\section{Phase 6: Production Hardening}
\textit{Objective: Secure and optimize the system for enterprise deployment.}
\begin{itemize}[label=$\bullet$, leftmargin=*]
    \item \textbf{Security:} Implement authentication, Role-Based Access Control (RBAC), and strict tenant-level query isolation across all endpoints.
    \item \textbf{Background Processing:} Extract heavy reconciliation jobs into a background worker queue (e.g., using BullMQ).
    \item \textbf{Observability:} Add structured logging (e.g., Pino) and track metrics for latency and matching performance.
    \item \textbf{Testing:} Write comprehensive unit tests for financial calculations and integration tests for the reconciliation loop.
    \item \textbf{Audit Logging:} Ensure every automated and manual decision writes an immutable entry to the \texttt{audit\_logs} table.
\end{itemize}

\end{document}
